\chapter{Low Blood Sugar (Hypoglycemia)}

\section*{Definition}
Hypoglycemia is a condition of abnormally low blood glucose levels, typically below 70 mg/dL (3.9 mmol/L). It is most commonly associated with diabetes treatment but can occur in various other conditions.

\section*{What Happens in Your Body}
The brain relies primarily on glucose for energy and cannot store significant amounts. When blood glucose falls below normal, the body releases counter-regulatory hormones (glucagon, epinephrine, cortisol, growth hormone) to raise glucose levels. The autonomic response causes adrenergic symptoms, while neuroglycopenia causes cognitive dysfunction. If prolonged, severe hypoglycemia can lead to seizures, coma, and brain damage.

\section*{Causes}
\begin{itemize}
\item Diabetes medications: insulin, sulfonylureas, meglitinides
\item Skipping or delaying meals
\item Increased physical activity without adjusting medication
\item Excessive alcohol intake (inhibits gluconeogenesis)
\item Critical illness: sepsis, liver failure, kidney failure
\item Insulinoma (insulin-secreting pancreatic tumor)
\item Hormonal deficiencies: adrenal insufficiency, hypopituitarism
\item Reactive hypoglycemia (after gastric surgery or in prediabetes)
\item Medications: beta-blockers, salicylates, quinolones
\item Factitious hypoglycemia (self-administration of insulin)
\end{itemize}

\section*{When to See a Doctor}
\begin{itemize}
\item Autonomic symptoms: sweating, tremor, palpitations, anxiety, hunger
\item Neuroglycopenic symptoms: confusion, drowsiness, difficulty speaking, weakness
\item Blurred vision or double vision
\item Headache
\item Irritability or mood changes
\item Seizures
\item Loss of consciousness or coma
\item Whipple triad: low blood glucose, symptoms, relief with glucose
\end{itemize}

\section*{Related Symptoms}
\begin{itemize}
\item Diabetes mellitus (type 1 and type 2)
\item Fatigue
\item Dizziness
\item Palpitations
\item Sweating
\item Seizures
\item Coma
\end{itemize}

\section*{Self-Care / Home Management}
\begin{itemize}
\item Check blood glucose if possible
\item Consume 15 grams of fast-acting glucose (glucose tablets, fruit juice, regular soda, candy)
\item Recheck blood glucose after 15 minutes
\item Repeat treatment if still below 70 mg/dL
\item Eat a protein-containing snack after recovery to prevent recurrence
\item If unconscious or unable to swallow, emergency glucagon is needed
\item Carry medical identification indicating diabetes
\end{itemize}

\section*{How Your Doctor Will Evaluate You}
\begin{itemize}
\item Capillary blood glucose measurement (fingerstick)
\item Detailed history of medication timing, meals, and symptoms
\item Fasting blood glucose, insulin, and C-peptide levels
\item 72-hour fasting test for suspected insulinoma
\item Mixed-meal tolerance test for reactive hypoglycemia
\item HbA1c to assess overall glycemic control
\end{itemize}

\section*{Treatment Options}
\begin{itemize}
\item Immediate glucose administration (oral for conscious, IV glucose for unconscious)
\item Glucagon injection (intramuscular or subcutaneous) for severe hypoglycemia
\item Adjust diabetes medications to prevent recurrence
\item Treat underlying cause (insulinoma resection, adrenal replacement)
\item Education on hypoglycemia recognition and management
\item Continuous glucose monitor for high-risk people
\end{itemize}

\section*{References}
\begin{thebibliography}{99}
\bibitem{ref1} Cryer PE, Axelrod L, Grossman AB, et al. "Evaluation and management of adult hypoglycemic disorders: an Endocrine Society clinical practice guideline." \textit{Journal of Clinical Endocrinology \& Metabolism}, 2009;94(3):709--728.
\bibitem{ref2} Seaquist ER, Anderson J, Childs B, et al. "Hypoglycemia and diabetes: a report of a workgroup of the American Diabetes Association and The Endocrine Society." \textit{Diabetes Care}, 2013;36(5):1384--1395.
\bibitem{ref3} Cryer PE. "Mechanisms of hypoglycemia-associated autonomic failure in diabetes." \textit{New England Journal of Medicine}, 2013;369(4):362--372.
\bibitem{ref4} Amiel SA, Dixon T, Mann R, et al. "Hypoglycaemia in type 2 diabetes." \textit{Diabetic Medicine}, 2008;25(3):245--254.
\bibitem{ref5} Cryer PE. "The prevention and correction of hypoglycemia." \textit{Physiological Reviews}, 2013;93(1):1--38.
\bibitem{ref6} Service FJ. "Hypoglycemic disorders." \textit{New England Journal of Medicine}, 1995;332(17):1144--1152.

\end{thebibliography}

\clearpage
